% Section 7 (group D): the common action W1/W2 and the Big Bang exploration.
% Sources: Companion rounds v4.9.0-r81--r86 (parts/v490-post-release-layers.tex) and
% v4.10.0-r01--r05 (parts/v4100-post-release-layers.tex).
\section{A common action W1/W2}
\label{sec:act}

\Conditional\ Everything in this section holds under decisions the author took between 4 and 6~November 2026, most of
them listed explicitly below. The derivations of the field equations are checked computations for the action as stated;
the action itself, its matter content and the way the tachyon sector enters it are choices. The section reports the
action programme as it stands in the Companion \cite{QSC}: the field equations follow consistently, the structure
supplies mechanisms, and it supplies no scale of its own. By an author decision of 6~November 2026 the programme is a
separate branch that is \emph{frozen}, without further repair decisions (\rnd{r05}, decision r05.3(iii)). Before that
decision about twelve author decisions had been taken since \rndold{r81}, most of them after a consequence did not fit;
each kept a consequence in range, and none produced a confirmed new prediction (\rnd{r05}, remark r05.2(c)). We state
this at the outset because it governs how the results below should be weighed.

\subsection{Decisions and the action}
\label{sec:act-decisions}

The action was set up in \rndold{r81} after the author fixed five points:
\begin{enumerate}
\item the coframe (tetrad) $e^a{}_\mu$ is an independent field; whether solutions satisfy the closure relation of the
  Companion is to be tested later, not assumed;
\item the main line is a real action with Dirac adjoint and symmetric kinetic term in the matter frame M; the complex
  reading (stationarity of $\mathrm{Re}\,S$ and $\mathrm{Im}\,S$) is computed alongside;
\item gravity is Einstein--Cartan gravity, without a cosmological term;
\item matter is stored in the Dixon carrier $\mathbb C\otimes\mathbb H\otimes\mathbb O$ \cite{Dixon1994} with the $U(1)$
  generated by $J=iL_{e_7}$, to be extended later by colour and the weak $SU(2)$;
\item the closure projector of the Companion is left open.
\end{enumerate}
The fields are the tetrad $e$, an independent Lorentz connection $\omega$ (so torsion is allowed), a $U(1)_J$ gauge field
$A$ with field strength $F$, and the sixteen left-handed Weyl states $\Psi$ of one generation including $\nu_R$, on which
the Lorentz group acts by left multiplication on $\mathbb H$ (\rndold{r81}, Theoremlet r81.1). The action W1 \emph{as
corrected} in \rnd{r01} reads, in signature $(+,-,-,-)$,
\begin{equation}
S_{\mathrm{W1}}=\int e\,\Bigl[-\frac1{2\kappa}R(e,\omega)
  +\frac12\bigl(\Psi^\dagger i\bar\sigma^\mu D_\mu\Psi+\mathrm{c.c.}\bigr)-\frac14F_{\mu\nu}F^{\mu\nu}\Bigr],
\label{eq:act-W1}
\end{equation}
with $e=\det e^a{}_\mu$ and
$D_\mu\psi=\partial_\mu\psi+\tfrac12\omega_{\mu ab}\Sigma^{ab}\psi+iqJA_\mu\psi$,
$\Sigma^{ab}=\tfrac14(\sigma^a\bar\sigma^b-\sigma^b\bar\sigma^a)$. As first written in \rndold{r81} the gravitational term
carried the opposite sign, $+\frac1{2\kappa}R$; Section~\ref{sec:act-sign} explains why it was changed.

Two properties hold on the full carrier (\rndold{r81}, r81.1(a), (c)):
\begin{itemize}
\item $U(1)_J$ is anomaly-free: $\tr q=\tr q^3=0$ for $J$, for $iR_{e_7}$ and for $B-L$. On the chiral image alone $B-L$
  would be anomalous ($\tr q^3=-1.78$), which is one reason $\Psi$ is the whole carrier and not that image.
\item Without a scalar potential no term proportional to $g_{\mu\nu}$ with a \emph{constant} coefficient arises: W1
  produces no vacuum term.
\end{itemize}

\paragraph{W2.} The action W2 of \rndold{r82} adds a Higgs doublet with potential
$V=-\mu^2|\Phi|^2+\lambda|\Phi|^4$ and Yukawa couplings through the connectors of the carrier that carry Higgs quantum
numbers. The sign of $-\mu^2$ is \emph{read} as the tachyonic transition connector of r81.9(c) (for a counter-moving
matter state and its tachyon image the connector momentum has $q^2=-2m^2\sinh2\eta<0$); it is not derived. Alone, W2 has
the vacuum energy $V_{\min}=-\mu^4/(4\lambda)$, about $(100\,\mathrm{GeV})^4$, with the familiar excess of order
$10^{55}$ (\rndold{r82}, r82.1(a)). Figure~\ref{fig:act-scheme} summarises the fields, their couplings and the field
equations that follow from each variation.

% F8: scheme of the action W1/W2 as layered sector planes (2.5D), vector, readable in black and white.
% Content from Companion rounds v4.9.0-r81, r82, r84, r86 and v4.10.0-r01. No numbers beyond those in Section 7.
\begin{figure}[t]
\centering
\begin{tikzpicture}[font=\small,>={Stealth[length=2.2mm]},
  plane/.style={draw=black,line width=0.6pt},
  fld/.style={align=center,inner sep=1.5pt,font=\footnotesize},
  eqbox/.style={draw,rounded corners=2pt,align=left,inner sep=3pt,text width=4.1cm,fill=white},
  tag/.style={anchor=south east,font=\scriptsize\sffamily,inner sep=1pt},
  chk/.style={line width=0.9pt},                % derived and checked
  set/.style={line width=0.9pt,dashed},         % set by author decision
  unc/.style={line width=0.9pt,densely dotted}  % not derived in the Companion
]
% geometry of a sector plane: width W, depth offset (dx,dy); shadow offset (sx,sy) for the second reading
\def\W{7.0}\def\dx{1.5}\def\dy{0.95}\def\sx{0.5}\def\sy{0.38}
\newcommand{\actsectorplane}[3]{% #1 = y offset, #2 = fill gray level, #3 = name
  \draw[dashed,line width=0.45pt,fill=white] (\sx,#1+\sy) -- (\W+\sx,#1+\sy) -- (\W+\dx+\sx,#1+\dy+\sy) -- (\dx+\sx,#1+\dy+\sy) -- cycle;
  \fill[black!#2] (0,#1) -- (\W,#1) -- (\W+\dx,#1+\dy) -- (\dx,#1+\dy) -- cycle;
  \draw[plane] (0,#1) -- (\W,#1) -- (\W+\dx,#1+\dy) -- (\dx,#1+\dy) -- cycle;
  \draw[line width=1.4pt] (0,#1) -- (\W,#1);   % heavier front edge: depth cue
  \node[anchor=north west,font=\small\bfseries,inner sep=2pt] at (0,#1) {#3};}
% the three sectors, back planes lighter
\actsectorplane{0}{15}{Higgs sector (W2)}
\actsectorplane{2.5}{10}{fermion sector}
\actsectorplane{5.0}{5}{gravity sector}
\node[anchor=south west,align=left,font=\footnotesize] at (0,5.0+\dy+\sy+0.05)
  {dashed planes: the same fields in the second reading\\
   (tachyon part, \textsc{c}\,4.9-r84): weight $a-b$ on derivative\\
   terms, $a+b$ on potentials and masses};
% field contents at the plane centres
\node[fld] (G) at (0.5*\W+0.5*\dx,5.0+0.5*\dy) {tetrad $e^a{}_\mu$,\enspace connection $\omega=\tilde\omega+K$\\[1pt] $-\tfrac1{2\kappa}\,e\,R(e,\omega)$, torsion in $K$};
\node[fld] (F) at (0.5*\W+0.5*\dx,2.5+0.5*\dy) {$\Psi$: 16 Weyl states,\enspace $A$: $U(1)_J$\\[1pt] $\tfrac12(\Psi^\dagger i\bar\sigma^\mu D_\mu\Psi+\mathrm{c.c.})-\tfrac14F^2$};
\node[fld] (H) at (0.5*\W+0.5*\dx,0.5*\dy) {Higgs doublet $\Phi$\\[1pt] $V=-\mu^2|\Phi|^2+\lambda|\Phi|^4$};
% couplings between the planes
\draw[chk,<->] (2.8,2.5+\dy+0.05) -- node[right,align=left,font=\footnotesize]{$D_\mu$: spin\\connection} (2.8,4.95);
\draw[chk,<->] (5.9,2.5+\dy+0.05) -- node[right,align=left,font=\footnotesize]{contact term\\$+\tfrac{3}{16}\kappa\,j\cdot j$} (5.9,4.95);
\draw[set,<->] (4.3,\dy+0.05) -- node[right,align=left,font=\footnotesize]{Yukawa via the\\connectors} (4.3,2.45);
% variations and field equations (right column)
\begin{scope}[xshift=10.35cm]
\node[eqbox,anchor=north west] (E1) at (0,6.95) {$\delta\omega$:\enspace Cartan equation,\\ $K_{cab}=\tfrac\kappa4\varepsilon_{cabd}j^d$};
\node[eqbox,anchor=north west] (E2) at (0,5.55) {$\delta e$:\enspace $\tilde G_{\mu\nu}=\kappa\,T^{\rm eff}_{\mu\nu}$,\\ compatible with Bianchi};
\node[eqbox,anchor=north west] (E3) at (0,4.15) {$\delta\psi^\dagger$:\enspace Weyl equation with\\ $+\tfrac{3\kappa}{8}j_a\bar\sigma^a\psi$;\enspace $\partial_\mu j^\mu=0$};
\node[eqbox,anchor=north west] (E4) at (0,2.75) {$\delta A$:\enspace Maxwell equation};
\node[eqbox,anchor=north west] (E5) at (0,1.55) {$\delta\Phi$:\enspace Higgs equation\\ (not derived in the Companion)};
\end{scope}
\draw[chk,->] (\W+\dx-0.3,5.0+0.75*\dy) to[out=10,in=180] (E1.west);
\draw[chk,->] (\W+0.3,5.0+0.2*\dy) to[out=0,in=180] (E2.west);
\draw[chk,->] (\W+\dx-0.3,2.5+0.75*\dy) to[out=15,in=180] (E3.west);
\draw[chk,->] (\W+0.3,2.5+0.2*\dy) to[out=0,in=180] (E4.west);
\draw[unc,->] (\W+0.6,0.5*\dy) to[out=0,in=180] (E5.west);
% status tags above the boxes
\node[tag,text=st-est] at (E1.north east) {established};
\node[tag,text=st-est] at (E2.north east) {established};
\node[tag,text=st-est] at (E3.north east) {established};
\node[tag,text=st-open] at (E5.north east) {open};
% legend
\begin{scope}[yshift=-0.95cm]
\draw[chk] (0,0) -- (0.8,0); \node[anchor=west,font=\footnotesize] at (0.85,0) {derived and checked (H88--H90)};
\draw[set] (5.6,0) -- (6.4,0); \node[anchor=west,font=\footnotesize] at (6.45,0) {set by author decision};
\draw[unc] (10.4,0) -- (11.2,0); \node[anchor=west,font=\footnotesize] at (11.25,0) {not derived};
\end{scope}
\end{tikzpicture}
\caption{The action W1/W2 as three sector planes \Conditional. Gravity (tetrad and independent connection with torsion,
corrected sign of $R$), the fermions stored in the Dixon carrier with the gauge field of $U(1)_J$, and the Higgs doublet of
W2. Vertical arrows are couplings: the spin connection in $D_\mu$, the contact term left by eliminating torsion, and the
Yukawa couplings through the connectors, whose form and the sign of $-\mu^2$ are set, not derived. The right column shows
which field equation follows from which variation; the first three were derived independently in Wolfram Language and
in a jet computation (\rndold{r86}, \rnd{r01}); the Maxwell equation was imposed there, not tested separately. The dashed planes behind are the tachyon part as a second reading of the
same fields (\rndold{r84}). Line styles carry the status in black and white.}
\label{fig:act-scheme}
\end{figure}


\subsection{Field equations, derived independently}
\label{sec:act-equations}

The first derivations (\rndold{r81}) were Python and sympy scripts written in the same sessions. In \rndold{r86} a second
system, Wolfram Language, derived the field equations from their definitions (pre-registered tests H88, H89, not blind,
since the expected forms were known), and \rnd{r01} closed the last gap with the fully non-linear Einstein--Cartan
equation (H90, not blind). We state the results for the corrected action~\eqref{eq:act-W1}; Section~\ref{sec:act-sign}
lists the sign changes relative to the action as first written.

\begin{proposition}[Field equations of W1; \rndold{r86}, \rnd{r01}]
\label{prop:act-fieldeq}
For the action \eqref{eq:act-W1}, with $j^a=\sum\Psi^\dagger\bar\sigma^a\Psi$ the total current of all sixteen states:
\begin{enumerate}
\item[(a)] \emph{Generator.} $\Sigma^{ab\dagger}\bar\sigma^c+\bar\sigma^c\Sigma^{ab}=-(\eta^{bc}\bar\sigma^a-\eta^{ac}\bar\sigma^b)$, hence
  $D_\mu j^c=\partial_\mu j^c+\omega_\mu{}^c{}_b\,j^b$; the ordering $\tfrac14(\bar\sigma^a\sigma^b-\bar\sigma^b\sigma^a)$ is not covariant.
\item[(b)] \emph{Cartan equation.} The $\omega$-equation is solved by $\omega=\tilde\omega+K$, with $\tilde\omega$ the torsion-free
  connection of the tetrad and the totally antisymmetric contorsion
  \begin{equation}
  K_{cab}=\frac{\kappa}{4}\,\varepsilon_{cabd}\,j^d .
  \label{eq:act-contorsion}
  \end{equation}
\item[(c)] \emph{Elimination of torsion.} Inserting (b) gives the four-fermion contact term
  \begin{equation}
  \mathcal L_{4f}=+\frac{3}{16}\,\kappa\,j\cdot j ,
  \label{eq:act-4f}
  \end{equation}
  with the total current; the torsion is blind to the $U(1)$ charge. For a single Weyl spinor the term vanishes, since its
  current is null.
\item[(d)] \emph{Weyl equation.} Varying with respect to $\psi^\dagger$ gives
  $i\bar\sigma^\mu\partial_\mu\psi-qJA_\mu\bar\sigma^\mu\psi+\tfrac{3\kappa}{8}\,j_a\bar\sigma^a\psi=0$ (written in the local
  form of the check).
\item[(e)] \emph{Current conservation.} The operator $M=qJA_\mu\bar\sigma^\mu-\tfrac{3\kappa}{8}j_a\bar\sigma^a$ is hermitian,
  hence $\partial_\mu j^\mu=0$ on solutions (Ward identity).
\item[(f)] \emph{Einstein--Cartan equation.} The tetrad variation at $\omega=\tilde\omega+K$ is symmetric on the field
  equations and gives, fully non-linearly,
  \begin{equation}
  \tilde G_{\mu\nu}=\kappa\,T^{\mathrm{eff}}_{\mu\nu},\qquad
  T^{\mathrm{eff}}_{\mu\nu}=T^{\mathrm{Bel}}_{\mu\nu}-g_{\mu\nu}\mathcal L_W+T^{\mathrm{Maxwell}}_{\mu\nu}
    -\frac{3\kappa}{16}\,g_{\mu\nu}\,j\cdot j ,
  \label{eq:act-einstein}
  \end{equation}
  where $\tilde G$ is the Einstein tensor of the Levi-Civita connection, $T^{\mathrm{Bel}}$ the Belinfante tensor and
  $\mathcal L_W$ the Weyl Lagrangian.
\item[(g)] \emph{Compatibility.} $\tilde\nabla^\mu T^{\mathrm{eff}}_{\mu\nu}=0$ on the field equations, consistent with the Bianchi
  identity $\tilde\nabla^\mu\tilde G_{\mu\nu}=0$.
\end{enumerate}
\end{proposition}

\begin{proof}[How it was checked]
(a), (c)--(e) were derived symbolically in Wolfram Language from the definitions: all $64$ components of the generator
identity, all $24$ contorsion components for two Weyl spinors (\rndold{r86}, r86.2). (b), (f), (g) were computed in
\rnd{r01} at a generic point of three random curved backgrounds (tetrad a random non-diagonal polynomial of degree three),
with all fields carried as jets to third order: three Weyl fields with charges $J=(+1,-1,+1)$, a real potential $A$ and
the tetrad, the spinor jets solving the Weyl equation with gauge and four-fermion terms at the point and to first
derivative. The residuals are $\le2\times10^{-14}$ for the Cartan equation, $10^{-15}$ for the Einstein equation and
$4\times10^{-15}$ (relative) for the divergence. Negative controls fail as they should: without the four-fermion term the
residuals rise to $0.12$--$0.19$ (Einstein) and $0.25$--$0.81$ (divergence); with spinors off the field equations to
$0.16$--$0.19$ and $0.92$--$1.05$. In addition, the first-order tetrad variation around flat space, with the spin
connection solved from $de+\omega\wedge e=0$ rather than inserted, equals minus the Belinfante tensor on a solution of
two plane waves, exactly symmetric, to $7\times10^{-15}$ (\rndold{r86}, r86.3(c)).
\end{proof}

These are numerical-exact checks at generic points, not a general symbolic proof of~(f); the round records them as such,
together with two errors found in the checking code (a conjugate transpose that also conjugated real symbols, and extra
$\eta$ factors in $\omega_{\mu ab}\Sigma^{ab}$). After their correction everything agreed; the disagreement lay in the
checking code, not in W1 (\rndold{r86}, disclosure to r86.3). The modulus $3/16$ of~\eqref{eq:act-4f} is the
Hehl--Datta value \cite{D:HehlDatta}, which the control of \rndold{r81} reproduced for a Dirac field.

\subsection{The sign of the gravitational action}
\label{sec:act-sign}

The non-linear computation of \rnd{r01} carried the sign of the gravitational term as a parameter,
$S_{\mathrm{grav}}=\frac{\sigma_g}{2\kappa}\int e\,R$, and found $\tilde G=-\sigma_g\kappa T^{\mathrm{eff}}$. For W1 as written
in \rndold{r81} ($\sigma_g=+1$) the factor is $-1$ to $4\times10^{-15}$. A weak-field computation in Wolfram fixes which
sign attracts: for $g=\mathrm{diag}\bigl(1+2\Phi,-(1-2\Phi)\mathbf 1\bigr)$ in the same curvature convention one finds
$G_{00}=2\nabla^2\Phi$. With $T_{00}=\rho>0$, $\tilde G=+\kappa T$ gives $\nabla^2\Phi=\kappa\rho/2=4\pi G\rho$ for
$\kappa=8\pi G$, which is Newtonian attraction, while $\tilde G=-\kappa T$ is repulsion. \Refuted\ W1 as first written
repels positive energy. The earlier Hehl--Datta control had compared only the modulus $3/16$, not the sign.

On 6~November 2026 the author corrected the action to~\eqref{eq:act-W1} (\rnd{r01}, decision r01.2). Moduli and forms of
all derived terms stay; the signs that change are (r01.3):
\begin{itemize}
\item the contact term becomes $\mathcal L_{4f}=+\tfrac3{16}\kappa\,j\cdot j$ and the Weyl-equation term
  $+\tfrac{3\kappa}{8}j_a\bar\sigma^a\psi$ (both were negative), the contorsion $K_{cab}=+\tfrac\kappa4\varepsilon_{cabd}j^d$;
\item the four-fermion stress becomes $T^{(4f)}_{\mu\nu}=-\tfrac{3\kappa}{16}\,j^2g_{\mu\nu}$ (Section~\ref{sec:act-stiff});
\item \emph{the torsion pairing channel becomes repulsive.} By the Fierz identity
  $\bar\sigma^\mu_{ab}\bar\sigma_{\mu\,cd}=2\varepsilon_{ac}\varepsilon_{bd}$ one has $j_1\cdot j_2=2|\Phi|^2$ for the Lorentz-scalar pair
  $\Phi=\psi_2^{\mathsf T}\varepsilon\psi_1$ (\rndold{r82}, r82.4). With the physical sign the interaction contains
  $H_{\mathrm{int}}\supset+\tfrac34\kappa\,\Phi^\dagger\Phi$; no condensate forms through it. The Planck-scale $\nu^c$
  Majorana condensate and the seesaw estimate that \rndold{r82} had drawn from the attractive sign are withdrawn.
\end{itemize}
Tests that do not involve the sign of the gravitational coupling are unaffected; \rnd{r01} lists them (r01.3(e)).

\subsection{A homogeneous fermion state is stiff matter, not a vacuum term}
\label{sec:act-stiff}

Since $T^{(4f)}_{\mu\nu}=-g_{\mu\nu}\mathcal L_{4f}$, the contact term taken \emph{on its own} has the algebraic form of a
cosmological term. In signature $(+,-,-,-)$ and a local rest frame,
\begin{equation}
\rho_{4f}=T^{(4f)}_{00}=-\mathcal L_{4f}=-\frac{3\kappa}{16}\,j^2,\qquad
p_{4f}=T^{(4f)}_{ii}=+\mathcal L_{4f}=-\rho_{4f},
\label{eq:act-4f-stress}
\end{equation}
a negative vacuum-like term, but with a height $\propto j^2$ that is not constant (\rndold{r86}, r86.3(a); sign after
\rnd{r01}). It cannot be separated physically from the kinetic part. For the whole tensor the answer is different:
two Weyl fields of opposite spin with a homogeneous phase solve the Weyl equation of
Proposition~\ref{prop:act-fieldeq}(d) exactly, with total current $(n,0,0,0)$, and on this solution the whole tensor
$T_{\mu\nu}=T^{\mathrm{Bel}}_{\mu\nu}-g_{\mu\nu}\mathcal L$ has
\begin{equation}
\rho=p=-\frac{3\kappa}{16}\,n^2 ,
\label{eq:act-stiff}
\end{equation}
that is, stiff matter ($w=1$) with the negative spin energy density known from Einstein--Cartan theory (\rndold{r86},
r86.3(b); \rnd{r01}, r01.3(b)). \Established\ (for W1 as stated) The homogeneous fermion state is stiff matter, not
$\Lambda$; W1 still produces no vacuum term.

\subsection{The tachyon image of an action}
\label{sec:act-image}

The Companion's frames M and T are related by the tachyon map, which on rapidities is the continuation
$\eta\mapsto\eta+i\pi/2$ (Convention F; it is not time reversal, and on the field doublet it is the electromagnetic
duality, not a Wick rotation). For a single momentum this is clear; for an \emph{action} it had not been fixed. Round
\rndold{r83} compared three readings for a scalar with $L_M=\tfrac12(\chi_t^2-\chi_x^2)-\tfrac12m^2\chi^2$ against three
criteria fixed in advance (pre-registered test H77, not blind): K1 tachyonic dispersion, as the frame map requires;
K2 $L_T=-L_M$ for massless fields, as in \rndold{r81}; K3 a running mode of the image carries energy of the opposite sign.

\begin{center}
\small
\begin{tabular}{@{}l>{\raggedright\arraybackslash}p{0.36\textwidth}lccc@{}}
\toprule
reading & rule & dispersion & K1 & K2 & K3\\
\midrule
A & $L_T=-L_M$ (energy parity) & $\omega^2=k^2+m^2$ & no & yes & yes\\
B & kinematics continued, $X\mapsto-X$, constants and potentials kept & $\omega^2=k^2-m^2$ & yes & yes & yes\\
B$'$ & $m^2\mapsto-m^2$ with real kinematics & $\omega^2=k^2-m^2$ & yes & no & no\\
\bottomrule
\end{tabular}
\end{center}

\noindent Only reading B is compatible; reading A is a ghost copy with normal dispersion. In reading B a potential without
derivatives keeps its sign, so the tachyonic Higgs mass of M is a normal mass in the image and the image does not
condense. The cancellation of the Higgs vacuum energy at photon weights that \rndold{r82} had found in reading A
therefore does not hold and was withdrawn; the vacuum energy is $a\,V_{\min}$ (\rndold{r83}, correction to r82.1).

\paragraph{Second reading of the same field.} If the tachyon part were an independent field, the vacuum of matter plus
image would decay through graviton exchange, $0\to\varphi\varphi\chi\chi$, at a rate that is unbounded without a preferred
frame; with a frame cutoff the gamma-ray background bounds the cutoff to about $14\,\mathrm{MeV}$ (\rndold{r83}, H76, not
blind; the scaling is that of \cite{D:Cline}). On 5~November 2026 the author decided instead that the tachyon part is
\emph{the same field in a second reading} (\rndold{r84}). The map is the complete complex boost with rapidity $i\pi$:
\begin{equation}
S(\theta)=e^{-\theta\sigma_z/2},\quad S\bar\sigma^\mu S=\Lambda(\theta)^\mu{}_\nu\bar\sigma^\nu,\qquad
S(i\pi)=-i\sigma_z,\quad \Lambda(i\pi)=\mathrm{diag}(-1,1,1,-1),
\label{eq:act-boost}
\end{equation}
acting on the spinor index \emph{and} on the argument, $\psi_T(x)=S\,\psi(\Lambda x+\delta)$. With the adjoint mapped
holomorphically, $\psi^\dagger\mapsto\psi^\dagger S$, the image term is $-L_M$ in all four directions, and the weighted action
is $(a-b)S_M$: isotropic, stable for $a>b$, with opposite energy in the image (criterion K3). The hermitian alternative
$\psi^\dagger\mapsto\psi^\dagger S^\dagger$ gives $(a+b)S_M$ and violates K3 (H79, not blind). The author fixed the holomorphic
form; in the matter reading the Dirac adjoint stays. A first form that rotated only the spinor index gave an anisotropic
Weyl operator with factor $(a+b)/(a-b)=\gamma_L=1.8645$ at matter weights, a Lorentz violation of order one for fermions
that laboratory bounds exclude (H78; \cite{Altschul2010}).

For a scalar the weighted Lagrangian is
\begin{equation}
a\,L_M+b\,L_T=(a-b)\,\tfrac12X-(a+b)\,V ,
\label{eq:act-weights}
\end{equation}
independent of the offset $\delta$ in flat space. Terms with derivatives enter with $a-b$, mass terms (Dirac and
Majorana masses are invariant under the complete boost) with $a+b$ (\rndold{r85}, H82). Photon weights are $a=b$; matter
weights have $b/a=\bsym^2$, so that $(a-b)/(a+b)=1/\gamma_L$. Since there are no independent T quanta, the decay channel
of H76 does not exist. \Conditional\ As a second reading of the same field, formed by the complete complex boost with
holomorphic adjoint, the tachyon part is isotropic, stable for $a>b$ and free of vacuum decay. This supports the
construction, not the physical identity of matter and tachyon sector (thesis T1 stays a conjecture). In flat space the
second reading only renormalises the kinetic term against the potential; the vacuum energy is cancelled in no variant.

\subsection{What the action does not supply}
\label{sec:act-noscale}

\Open\ The only scale W1 and W2 bring with them is the Planck scale; $\mu$ and $\lambda$ of W2 are inputs
(\rndold{r82}, consequence). A $\Lambda$-like term can arise from tachyonic modes of both norm signs, but its height and
sign are set by a free weight ratio (\rndold{r81}, H72); with the second reading no tachyonic branch remains, and every
route examined in \rnd{r03} and \rnd{r04} ends with a vacuum term that is a free multiple of $V_{\min}$. The aim of the
programme, a vacuum term that is not free, is not met (\rnd{r05}, r05.2(c)); Section~\ref{sec:lop5} states the corresponding negative
result for $\Lambda$ in the full class of constructions of the Companion. Table~\ref{tab:act-status} collects the status.

\begin{table}[ht]
\centering
\small
\begin{tabularx}{\textwidth}{@{}>{\raggedright\arraybackslash}X>{\raggedright\arraybackslash}p{0.43\textwidth}l@{}}
\toprule
statement & status & source\\
\midrule
action W1, $U(1)_J$ anomaly-free on the full carrier & \Conditional\ set up under author decisions & \rndold{r81}\\
field equations (generator, torsion, Weyl, Ward) & \Established\ for W1, derived in Wolfram (H88) & \rndold{r86}\\
non-linear Einstein--Cartan equation, Bianchi & \Established\ for W1, derived and compatible (H90) & \rnd{r01}\\
sign of the gravitational action & \Refuted\ as first written (repulsion); corrected by author decision & \rnd{r01}\\
torsion pairing channel & repulsive after the correction; Planck-scale condensate withdrawn & \rnd{r01}\\
homogeneous fermion state & stiff matter, $\rho=p<0$; not a vacuum term & \rndold{r86}, \rnd{r01}\\
tachyon image of an action & reading B (kinematics continued); A excluded (H77) & \rndold{r83}\\
tachyon part as second reading & \Conditional\ isotropic, stable, no vacuum decay (H79) & \rndold{r84}\\
Higgs vacuum energy with image & no cancellation, $a\,V_{\min}$ & \rndold{r83}\\
own scale, value of $\Lambda$ & \Open\ none supplied; programme frozen & \rnd{r05}\\
\bottomrule
\end{tabularx}
\caption{Status of the action programme W1/W2. ``Established'' refers to the derivation for the stated action, not to
the action as physics.}
\label{tab:act-status}
\end{table}

\subsection{The Big Bang: an exploration}
\label{sec:act-bang}

\Exploration\ The second reading suggests a picture of the Big Bang, which the Companion has explored; most tests on it
were pre-registered but not blind. By a further author decision, the domain choice G2 of \rndold{r85}, the matter reading
lives after the Big Bang and the tachyon reading on its image before it, joined only at $\tau=0$. Structurally this is the
two-sheet picture of the CPT-symmetric universe of Boyle, Finn and Turok \cite{D:BFT}; the resemblance was noticed
afterwards.

\paragraph{What is shown.}
\begin{itemize}
\item \emph{A CPT-like transition with an admixture fixed by $\bsym$} (\rndold{r85}). For conformal fields the weight jumps
  from $-b$ to $a$. In the sharp limit a positive-norm mode before the bang becomes
  $C\,e^{-i\omega\tau}+D\,e^{+i\omega\tau}$ with $C=(a+b)/(2a)$, $D=(a-b)/(2a)$; normalised, the Bogoliubov ratio is
  $|\beta/\alpha|=1/\gamma_L=0.536$ and the occupation $n=(1-\bsym^2)^2/(4\bsym^2)=0.404$ per mode, for every $k$; at photon
  weights $n=0$. With a smooth transition of duration $\Delta$, taken around zero in the complex plane, every mode of the
  tachyon reading is mapped to a mode of negative frequency after the bang (a CPT assignment), and the admixture is cut
  off exponentially for $k\Delta\gg1$ (H86; checked by a second system in \rndold{r86}). $\Delta$ is a free scale.
\item \emph{Fermions pass without admixture} (\rnd{r03}). The weighted Dirac or Weyl equation is solved exactly by
  $\psi=c^{-1/2}\chi$ with $\chi$ free, for every profile $c(\tau)$; positive and negative frequencies do not mix, and the
  transfer factor is $\pm i\bsym$, unitary in the weighted norm. A first-order equation lets the weight be scaled away; the
  second-order equations of bosons do not. On an expanding background this holds for $k\gg m\,a(\tau)$ at the transition.
\item \emph{Bosons are norm-preserving only far from $1/\Delta$} (\rnd{r03}, H94). For $k\Delta\ll1$ they pass as a
  Bogoliubov transformation, for $k\Delta\gg1$ they are CPT-mirrored; in between the transition is not norm-preserving.
  The author set $\Delta\lesssim\ell_P$, which moves that range to the Planck scale.
\item \emph{No $\Lambda$ from the tachyon sector after the bang} (\rnd{r03}, \rnd{r04}). Under G2 only the matter reading
  exists after the bang, so neither the tension of a matter--tachyon mixture nor a tachyonic or ghost condensate survives
  it. A vacuum term after the bang can only be $a\,V_{\min}$ of the matter reading, with a free value. Admitting the
  tachyonic branch as degrees of freedom of its own did not help either (H95 refuted), and that admission was withdrawn.
  $\Lambda$ stays an input.
\end{itemize}

\paragraph{What could still be.} \Open\ Whether the transition has a physical duration $\Delta$ and what fixes it; what
the two senses of the complex rotation mean (they are exactly conjugate and give an asymmetry with fixed product); how
tensor modes see the transition, which an observable trace such as a gravitational-wave background would need; and the
energy in bosons created near $k\sim1/\Delta$, which sets a Planck-hot, non-thermal start without fermions and an open
problem for $\Delta N_{\mathrm{eff}}$ (\rnd{r04}, \rnd{r05}). None of this is derived. The picture is a consistent reading
of the second reading under G2, not a result about the early universe.
